\documentclass[handout]{beamer}
\mode<presentation>
\usepackage{xcolor}
\usepackage[toc]{multitoc}
\usepackage{color}
%\usepackage{enumitem}
\usetheme[secheader]{Boadilla}
\usecolortheme{whale}
\setbeamertemplate{caption}[numbered]
\setbeamertemplate{bibliography item}[text]
\setbeamertemplate{navigation symbols}{}

\theoremstyle{definition}
\newtheorem{definisi}[theorem]{Definisi}
\theoremstyle{example}
\newtheorem{contoh}[theorem]{Contoh}

\renewcommand*{\multicolumntoc}{2}

%#make sure to change this part, since it is predefined
      %\defbeamertemplate*{footline}{infolines theme}
      \setbeamertemplate{footline}
        {
      \leavevmode%
      \hbox{%
      \begin{beamercolorbox}[wd=.333333\paperwidth,ht=2.25ex,dp=1ex,center]{author in head/foot}%
        \usebeamerfont{author in head/foot}\insertshortauthor%~~(\insertshortinstitute)
      \end{beamercolorbox}%
      \begin{beamercolorbox}[wd=.433333\paperwidth,ht=2.25ex,dp=1ex,center]{title in head/foot}%
        \usebeamerfont{title in head/foot}\insertshorttitle
      \end{beamercolorbox}%
      \begin{beamercolorbox}[wd=.233333\paperwidth,ht=2.25ex,dp=1ex,right]{date in head/foot}%
        \usebeamerfont{date in head/foot}\insertshortdate{}\hspace*{2em} \insertframenumber{} / \inserttotalframenumber\hspace*{2ex} 
      \end{beamercolorbox}}%
      \vskip0pt%
    }

%\AtBeginSection[] {
%    \begin{frame}<beamer>{Daftar Isi}
%    \frametitle{Daftar Isi}
		%\tableofcontents[currentsection,currentsubsection,sectionstyle=show/hide,subsectionstyle=show/show/hide,subsubsectionstyle=show/show/show/hide]
   % \end{frame}
%}


\subtitle[Logika Matematika]{Logika Matematika}
\title[Logika Predikat (Kalkulus Predikat)]{Logika Predikat (Kalkulus Predikat)}
\author{Cecep Suwanda, S.Si., M.Kom.}
\date{Juni 9, 2025}
\institute{Universitas Bale Bandung}

\begin{document}



\maketitle

%\begin{frame}
%\titlepage
%\end{frame}

\section*{Daftar Isi}
\begin{frame}[shrink]{Daftar Isi}
\frametitle{Daftar Isi}
\tableofcontents
\end{frame}

\section{Pendahuluan}
\begin{frame}
\frametitle{Pendahuluan}

Kalimat pada kalkulus proposisi tidak dapat menjelaskan \textit{konsep objek} dan \textit{relasi antar objek}.\\
Contoh
\begin{center}
	\color{blue}
	\textit{Batuan di Mars berwarna putih}\\
atau\\
\textit{Batuan di Mars tidak berwarna putih}\\
\end{center}
Dengan aturan kalkulus proposisi, pernyataan tersebut dapat dibuat menjadi skema kalimat\\
\begin{center}
	\textit{(p or not p)}\\
\end{center}
dan selanjutnya dapat ditentukan nilai kebenarannya

\end{frame}

\begin{frame}
	\frametitle{Pendahuluan}	
	Jika ada pernyataan lain,\\
	\begin{center}
	\color{blue}	
		\textit{ Ada batuan di Mars berwarna putih}\\
	atau\\
	\textit{Semua batuan di Mars berwarna putih}\\
	\end{center}
	maka pernyataan di atas tidak dapat dibentuk menjadi skema kalimat kalkulus	proposisi.\\
	Hal ini disebabkan karena pernyataan tersebut mengandung {\color{blue} kuantisasi dari objek}.\\
	Oleh karena itu dibutuhkan bahasa baru yang mengenal adanya konsep objek dan relasi antar objek, yaitu menggunakan {\color{blue} Kalkulus Predikat}.	
\end{frame}

\begin{frame}
	\frametitle{Pendahuluan}
	Dengan kalkulus predikat maka pernyataan tersebut diubah menjadi :\\
	\begin{center}
		\color{blue}
		(for some x) (p(x) and q(x))\\
		or\\
		(for all x)(if p(x) then q(x))\\
	\end{center}
	dimana :\\
	p(x) = x adalah batuan di Mars \\
	q(x) = x adalah batuan berwarna putih\\
	"for some x" disebut kuantifier (simbol : $\exists x$) \\
	"for all x" disebut kuantifier (simbol : $\forall x$) \\
\end{frame}

\begin{frame}
	\frametitle{Pendahuluan}
	Pada dasarnya Kalkulus Predikat merupakan perluasan	dari Kalkulus Proposisi dimana Kalkulus Predikat mengatasi kelemahan pada kalkulus proposisi dengan menambahkan representasi
	\begin{enumerate}
		\item Objek yang memiliki sifat tertentu
	    \item Relasi antar objek
    \end{enumerate}	    
\end{frame}

\section{Definisi Simbol}
\begin{frame}
	\frametitle{Definisi Simbol}
	Kalimat dalam kalkulus predikat dibuat dari simbol-simbol berikut.\\
	\begin{enumerate}
		\item Simbol Kebenaran : true dan false
	 	\item Simbol Konstanta : $a, b, c, a_1, b_1, ...$
	 	\item Simbol variabel : $x, y, z, x_1, x_2, ...$
	 	\item Simbol fungsi : $f, g, h, g_1, f_1, h_1, ...$ \\
	Setiap simbol fungsi mempunyai arity yang menyatakan banyaknya parameter/argumen yang harus dipenuhi.
	 	\item Simbol Predikat (menyatakan relasi): $p, q, r, s, p_1, q_1,  ...$\\
	Setiap simbol predikat juga memiliki arity
    \end{enumerate}	 
\end{frame}	

\section{Definisi Term}
\begin{frame}
	\frametitle{Definisi Term}
	Term adalah sebuah ekspresi yang menyatakan objek.\\
	Term dibangun berdasarkan aturan-aturan sebagai berikut.\\
	\begin{enumerate}
		\item Semua konstanta adalah term
		\item Semua variabel adalah term
		\item Jika $t_1, t_2, ..., t_n$ adalah term ($n \geq 1$) dan f adalah fungsi dengan arity = n,
	maka fungsi f($t_1,t_2, ..., t_n$) adalah term
		\item Jika A adalah kalimat, sedang s dan t adalah term, maka kondisional if A
	then s else t adalah term
    \end{enumerate}
\end{frame}	

\begin{frame}
	\frametitle{Definisi Term}
    Contoh :\\
    \begin{enumerate} 
    	\item  f(a,x) adalah term, karena\\
    a adalah simbol konstanta, dan semua konstanta adalah term,\\
    x adalah simbol variabel, dan semua variabel adalah term,\\
    f adalah simbol fungsi dan semua fungsi adalah term
    	\item  g(x, f(a,x)) adalah term, karena\\
    a adalah simbol konstanta, dan semua konstanta adalah term,\\
    x adalah simbol variabel, dan semua variabel adalah term,\\
    f dan g adalah simbol fungsi dan semua fungsi adalah term\\
    \end{enumerate}
\end{frame}	

\section{Definisi Proposisi}
\begin{frame}
	\frametitle{Definisi Proposisi}
   Proposisi digunakan untuk merepresentasikan relasi antar objek\\
   Proposisi dibangun berdasarkan aturan sebagai berikut :\\
   \begin{enumerate} 
   	\item  Simbol kebenaran adalah proposisi
    \item  Jika $t_1, t_2, ..., t_n$ adalah term dan p adalah simbol predikat dengan n-ary maka p ($t_1,t_2, ..., t_n$) adalah proposisi
   \end{enumerate}
   Contoh :\\
   p (a, x, f (a,x)) adalah proposisi, karena\\
   a adalah simbol konstanta dan\\
   x adalah simbol variabel, dan\\
   f adalah simbol fungsi, dan\\   
   semua konstanta, variabel, dan fungsi adalah term dan\\
   p adalah simbol predikat 3-ary
\end{frame}		

\section{Definisi Kalimat}
\begin{frame}
	\frametitle{Definisi Kalimat}
      Kalimat dalam kalkulus predikat dibangun dengan aturan,\\
		\begin{enumerate} 
			\item  Setiap proposisi adalah kalimat,
		    \item  Jika A, B, C adalah kalimat maka
		          \begin{enumerate} 
			         \item \textbf{Negasi} (not A) adalah kalimat
					 \item \textbf{Konjungsi} A dengan B: (A and B) adalah kalimat
			         \item \textbf{Disjungsi} A dengan B : (A or B) adalah kalimat
			         \item \textbf{Implikasi} (If A then B) adalah kalimat
			         \item \textbf{Ekivalensi} A dan B (A if and only if B) adalah kalimat
			         \item \textbf{Kondisional} if A then B else C adalah kalimat.
			     \end{enumerate}   
		    \item  Jika A adalah kalimat dan x adalah variabel maka,
		          \begin{enumerate} 
		          	\item (\textbf{For all x}) A adalah kalimat
		          	\item (\textbf{For some x}) A adalah kalimat
		         \end{enumerate} 	
	    \end{enumerate}
\end{frame}

\begin{frame}
	\frametitle{Definisi Kalimat}
	Contoh :\\
	\begin{enumerate} 
		\item  if (for all x) p(a, b, x) then g (y) else f(a, y) adalah term, karena
			\begin{enumerate} 
				\item  a dan b adalah simbol konstanta,
				\item  x dan y adalah simbol variabel,
				\item  f dan g adalah simbol fungsi, dan
				\item  Semua konstanta, variabel dan fungsi adalah term.
				\item  p adalah simbol predikat.
				\item  (for all x) p(a, b, x) adalah kalimat, g (y) dan f(a, y) adalah term, maka kondisional if (for all x) p(a, b, x) then g (y) else f(a, y) adalah term
		     \end{enumerate}
		     
		\item  if (all for x) p(a, b, x) then (for some y) q(y) else not p(a,b,c) adalah kalimat
		 \begin{enumerate} 
		 	\item  a dan b adalah simbol konstanta,
		 	\item  x dan y adalah simbol variabel,
		 	\item  Semua konstanta dan variabel adalah term,
		 	\item  p dan q adalah simbol predikat,
		 	\item  (for all x) p(a, b, x) dan (for some y) q(y) adalah kalimat, 
		 	maka \\kondisional if (for all x) p(a, b, x) then (for some y) q(y) else \\not p(a, b, c) adalah kalimat 
		 \end{enumerate}    
    \end{enumerate}
	
\end{frame}	

\section{Definisi Ekspresi}
\begin{frame}
	\frametitle{Definisi Ekspresi}
     Suatu ekspresi dalam kalkulus predikat dapat berupa kalimat atau term.\\
     Contoh :\\
	  \begin{tabular}{ l l }	
		x        &         merupakan ekspresi\\
		f(x,y)   &         merupakan ekspresi\\
		(for some x) p(x) & merupakan ekspresi
	 \end{tabular}	
\end{frame}

\section{Definisi Subterm,Subkalimat,Subekspresi}
\begin{frame}
	\frametitle{Definisi Subterm,Subkalimat,Subekspresi}
	 \begin{enumerate} 
	 	\item  Subterm dari term t atau dari kalimat A adalah setiap term antara yang digunakan untuk membangun t atau A
	 	\item  Subkalimat adalah setiap kalimat antara yang digunakan untuk membangun kalimat yang lebih luas
	 	\item  Subekspresi adalah subterm atau subkalimat yang terdapat	pada sebuah ekspresi
     \end{enumerate}
     Contoh :\\
     Sebutkan semua subterm dan subkalimat yang terdapat pada ekspresi berikut :\\
     \begin{center}
     	\textbf{E : if (for all x) q (x, f(a)) then f (a) else b}
     \end{center}
     Subterm : a, x, f(a), b, if (for all x) q (x, f(a)) then f (a) else b\\
     Subkalimat : q(x, f(a), (for all x) q(x,f(a))\\
     Semuanya merupakan SubEkspresi dari E\\
\end{frame}	

\section{Representasi Kalimat}
\begin{frame}
	\frametitle{Representasi Kalimat}
      Contoh representasi bahasa alami dalam Kalkulus Predikat\\
      \begin{enumerate} 
      	\item  {\color{red} Ada} apel berwarna merah\\
      ({\color{red} FOR SOME x}) (Apel(x) AND Merah(x))
        \item  {\color{red} Semua} apel berwarna merah\\
         ({\color{red}FOR ALL x}) ( IF Apel(x) THEN Merah(x))
        \item {\color{red}Setiap} orang mencintai {\color{blue}seseorang}\\
        ({\color{red}FOR ALL x}) ({\color{blue}FOR SOME y}) LOVES(x,y) 
        \item Ani dicintai {\color{red}banyak} orang\\
        ({\color{red}FOR ALL x}) LOVES(x, Ani)
  \end{enumerate}  
 \end{frame}
        
\begin{frame}
	\frametitle{Representasi Kalimat}
	 \begin{enumerate}       
        \item {\color{red}Semua} Apel berwarna merah terasa manis\\
        ({\color{red}FOR ALL x}) (IF apel(x) AND merah(x) THEN manis(x))\\
        ({\color{red}FOR ALL x}) (IF apel(x) THEN (IF merah(x) THEN manis(x)))
        \item {\color{red}Tidak semua} apel berwarna merah terasa manis\\
        NOT {[} ({\color{red}FOR ALL x}) (IF apel(x) AND merah(x) THEN manis(x)){]} \\
        {[}NOT ({\color{red}FOR ALL x}){]} {[}NOT (IF apel(x) AND merah(x) THEN manis(x)){]}\\
        ({\color{red} FOR SOME x}) (apel(x) AND merah(x) AND NOT manis(x))
    \end{enumerate}      
\end{frame}

\section{Latihan Soal}
\begin{frame}
    \frametitle{Latihan Soal}
    \begin{enumerate}       
    	\item {\color{red} Semua} Komunis itu tidak bertuhan
    	\item Tidak {\color{red}ada} gading yang tidak retak
    	\item {\color{red}Ada} gajah yang jantan dan {\color{red}ada} yang betina
    	\item {\color{red}Tidak semua} pegawai negeri itu manusia korup
    	\item {\color{red}Hanya polisilah} yang berwenang mengadakan penyidikan, kalau {\color{red}ada} orang yang melanggar hukum
    	\item {\color{red}Semua} orang komunis itu bukan pancasilais.\\
              {\color{red}Ada} orang komunis yang anggota tentara.\\
              Jadi, {\color{red}ada} anggota tentara yang bukan pancasilais
    	\item {\color{red}Barang siapa} meminjam barang orang lain dan tidak mengembalikannya adalah penipu. {\color{red}Ada} penipu yang begitu lihai, sehingga tidak ketahuan.
        Kalau {\color{red}orang} menipu dan itu tidak ketahuan, ia tidak dapat dihukum.
        Jadi {\color{red}ada} penipu yang tidak dapat dihukum. 
   \end{enumerate}     	
\end{frame}

\section{Jawaban Latihan Soal}
\begin{frame}
	\frametitle{Jawaban Latihan Soal}
	\begin{enumerate}
	   \item {\color{red} Semua} Komunis itu tidak bertuhan\\
	         {\color{red} $\forall$x} [IF Komunis(x) THEN NOT Bertuhan(x)]\\
	   \item Tidak {\color{red}ada} gading yang tidak retak\\
	         NOT  [({\color{red}$\exists$x}) (Gading(x) AND NOT Retak(x)]\\
	         {\color{red}$\forall$x} (NOT (Gading(x) AND NOT Retak(x)]\\
	         {\color{red}$\forall$x} (IF Gading(x) THEN Retak(x))
	   \item {\color{red}Ada} gajah yang jantan dan {\color{red}ada} yang betina\\
	   ({\color{red}$\exists$x})[ (Gajah(x) AND Jantan(x)) OR (Gajah(x) AND Betina(x))]\\
	   ({\color{red}$\exists$x})[ (Gajah(x) AND (Jantan(x)) OR Betina(x))]
	   \item {\color{red}Tidak semua} pegawai negeri itu manusia korup\\
	    NOT [{\color{red}$\forall$x} IF Manusia(x) AND Pegawai\_Negeri(x) THEN Korup(x)]\\
	   ({\color{red}$\exists$x}) [NOT ( NOT(Manusia(x) AND Pegawai\_Negeri(x) AND NOT Korup(x)]\\
	   ({\color{red}$\exists$x}) [Pegawai\_Negeri(x) AND Manusia(x) AND NOT Korup(x)]      
	\end{enumerate}
\end{frame}

\begin{frame}
	\frametitle{Jawaban Latihan Soal}
	\begin{enumerate}
		\addtocounter{enumi}{4}
	    \item {\color{red}Hanya polisilah} yang berwenang mengadakan penyidikan, kalau {\color{red}ada} orang yang melanggar hukum\\
	     {\color{red}$\exists$x}{\color{red}$\forall$y}[IF Orang(x) AND MelanggarHukum(x) THEN (IF Polisi(y) THEN Penyidikan(y,x))]
	    \item {\color{red}Semua} orang komunis itu bukan pancasilais.\\
	          {\color{red}Ada} orang komunis yang anggota tentara.\\
	          Jadi, {\color{red}ada} anggota tentara yang bukan pancasilais\\
	          {\color{red}$\forall$x}[IF Komunis(x) THEN NOT Pancasilais(x)];\\ 
	          {\color{red}$\exists$x}[Komunis(x) AND Tentara(x)];\\
	          {\color{red}$\exists$x}[Tentara(x) AND NOT Pancasilais(x)]  
     \end{enumerate}
\end{frame}

\begin{frame}
	\frametitle{Jawaban Latihan Soal}
	\begin{enumerate}
		\addtocounter{enumi}{6}
		\item {\color{red}Barang siapa} meminjam barang orang lain dan tidak mengembalikannya adalah penipu. {\color{red}Ada} penipu yang begitu lihai, sehingga tidak ketahuan.
		Kalau {\color{red}orang} menipu dan itu tidak ketahuan, ia tidak dapat dihukum.
		Jadi {\color{red}ada} penipu yang tidak dapat dihukum.\\
		\begin{enumerate}
		\item (FOR ALL x,y,z) IF Orang(x) AND Orang(z) AND Barang(y) AND Meminjam(x,y,z) 
		AND NOT Mengembalikkan(x,y,z) THEN Penipu(x)
		\item (FOR SOME $x_1$) Penipu($x_1$) AND Lihai($x_1$) AND NOT Ketahuan($x_1$)
		\item (FOR ALL x) IF Penipu(x) AND NOT Ketahuan(x) THEN NOT Hukum(x)
		\item (FOR SOME $x_1$) Penipu($x_1$) AND NOT Hukum($x_1$)
	    \end{enumerate} 			
\end{enumerate}
\end{frame}

\section{Variabel Bebas atau Terikat}
\begin{frame}
	\frametitle{Variabel Bebas atau Terikat}
    \begin{enumerate}
    	\item Suatu variabel dikatakan terikat dalam sebuah ekspresi jika sedikitnya ada satu kemunculan x terikat pada ekspresi tersebut.
    	\item Sebaliknya dikatakan variabel bebas jika sedikitnya ada satu kemunculan bebas dalam ekspresi tersebut.
    \end{enumerate}
     Contoh : \\ 
     \begin{center}
     	\textbf{(FOR ALL x)[p(x,y) AND (FOR SOME y) q(y,z)]}
     \end{center}
      x pada p(x,y) adalah terikat\\
      y pada p(x,y) adalah bebas\\
      y pada q(y,z) adalah terikat\\
      z pada q(y,z) adalah bebas 
     
\end{frame}	

\begin{frame}
	\frametitle{Variabel Bebas atau Terikat}
	Kemunculan variabel terikat dipengaruhi oleh kemunculan kuantifier yang paling dekat.\\
	Contoh :\\
	\begin{center}
		\textbf{(FOR ALL x)[p(x) OR (FOR SOME x)(FOR ALL y) r(x,y)]}
	\end{center}
	 variabel x pada p(x) dipengaruhi kuantifier FOR ALL x\\
	 variabel x pada r(x,y) dipengaruhi kuantifier FOR SOME x\\
	 Catatan :\\
	 Perbedaan antara variabel Bebas dan Variabel Terikat adalah :\\
	 Variabel Bebas, Nilainya diberikan oleh interpretasi\\
	 Variabel Terikat,Nilainya terbatas dari interpretasi yang diberikan
	
\end{frame}	

\section{Kalimat Tertutup}
\begin{frame}
	\frametitle{Kalimat Tertutup}
	Sebuah kalimat dikatakan tertutup jika tidak mempunyai kemunculan bebas dari variabel-variabelnya.\\
	Contoh :\\
	\begin{enumerate}
		\item  (FOR ALL x) (FOR SOME y) p(x, y)adalah kalimat tertutup.
		\item (FOR ALL x) p(x, y) bukan merupakan kalimat tertutup
	\end{enumerate}		
\end{frame}	

\section{Simbol Bebas}
\begin{frame}
	\frametitle{Simbol Bebas}
	Simbol bebas dari ekspresi A adalah :\\
	\begin{itemize}
		\item variabel-variabel bebas
		\item semua konstanta
		\item semua simbol fungsi
		\item semua simbol predikat		
	\end{itemize}
	dari ekspresi A

\end{frame}

\section{Interpretasi}
\begin{frame}
	\frametitle{Interpretasi}
	 Misal D adalah sebarang himpunan tak kosong, maka sebuah interpretasi I dalam domain D akan memberi nilai pada setiap simbol konstanta,	variabel bebas, fungsi dan predikat yang ada pada kalimat dengan aturan sebagai berikut :\\
	 \begin{itemize}
	 	\item Untuk setiap konstanta a, yaitu elemen $a_1$ dari D
	 	\item Untuk setiap variabel x, yaitu elemen $x_1$ dari D
	 	\item Untuk setiap simbol fungsi f dengan arity = n , yaitu :\\
	 	Fungsi $f_1$($d_1$, $d_2$, …, $d_n$) dimana argumen $d_1$, $d_2$, …, $d_n$ merupakan elemen dari D, dan nilai fungsi $f_1$($d_1$, $d_2$, …, $d_n$) merupakan anggota D
	 	\item Untuk setiap simbol predikat p dengan arity = n, yaitu relasi $p_1$($d_1$, $d_2$, …, $d_n$) dimana argumen $d_1$, $d_2$, …, $d_n$ merupakan elemen dari D dan nilai $p_1$($d_1$, $d_2$, …, $d_n$) adalah TRUE atau FALSE
	 \end{itemize}
	  Jadi untuk suatu ekspresi A, sebuah interpretasi I dikatakan interpretasi untuk A, jika I memberikan nilai kepada setiap simbol bebas dari A. 
	
\end{frame}	

\section{Arti Kalimat}
\begin{frame}
	\frametitle{Arti Kalimat}
	\begin{itemize}
		\item Arti kalimat ditentukan oleh interpretasi yang diberikan. Tetapi karena dalam kalkulus predikat mengandung pengertian objek, maka interpretasi dalam kalimat predikat harus juga mendefinisikan suatu \textbf{\underline{domain yaitu himpunan objek yang memberi arti pada term}}.
		\item Suatu interpretasi harus memberi nilai pada setiap simbol bebas pada kalimat tersebut.
	\end{itemize}
	
	
\end{frame}	

\begin{frame}
	\frametitle{Arti Kalimat}
	Misalkan ada kalimat tertutup :\\
	\begin{center}
		\textbf{A : IF (FOR ALL x)(FOR SOME y) p(x,y) THEN p(a,f(a))}
	\end{center}
	Interpretasi untuk kalimat A harus\\
	\begin{itemize}
		\item Mendefinisikan Domain
		\item Memberikan nilai untuk simbol bebas dalam hal ini : Konstanta a, Simbol fungsi f, Simbol p
	\end{itemize}
	Contoh :\\
	 Diberikan interpretasi I dengan Domain D adalah himpunan bilangan integer positif, dimana :\\
	 a = 0\\
	 p = relasi “lebih besar” yaitu : p(d1,d2)=(d1 $>$ d2)\\
	 f  = fungsi suksesor yaitu f(d) = d + 1\\
	 berdasarkan interpretasi I, kalimat tersebut dapat diartikan sebagai :\\
	 IF untuk setiap integer x Ada integer y sedemikian sehingga x $>$ y  THEN 0 $>$ 0 + 1
		
\end{frame}

\begin{frame}
	\frametitle{Arti Kalimat}
	Misalkan interpretasi J dengan domain bilangan integer positif, yang akan memberi nilai :\\
	a = 0\\
	p = relasi ketidaksamaan yaitu : p(d1,d2)=(d1$\neq$d2)\\
	f  = fungsi predesesor yaitu f(d) = d - 1\\
	Berdasarkan interpretasi J, kalimat tersebut dapat diartikan sebagai :\\
	IF untuk setiap integer x Ada integer y sedemikian sehingga  x $\neq$ y  THEN 0 $\neq$ 0 – 1
\end{frame}

\section{Latihan Soal}
\begin{frame}
	\frametitle{Latihan Soal}
	 Diberikan Ekspresi :\\
	 \begin{center}
	   \textbf{E = IF p(x,f(x)) THEN (FOR SOME y) p(a,y)}
	\end{center}
	\begin{enumerate}
		\item Misalkan I adalah interpretasi untuk E dengan Domain bilangan real; dimana\\ 
		a = $\sqrt{2}$\\
		x = $\pi$\\
		f = fungsi “dibagi 2” yaitu : f1(d1) = d1/2\\
		p  = relasi “lebih besar atau sama dengan” yaitu p(d1,d2) = (d1$\geq$d2)
		\item  Misalkan J adalah interpretasi untuk E dengan Domain semua orang; dimana\\ 
		a = Soeharto\\
		x = Soekarno\\
		f = fungsi “Ibu dari” yaitu : f1(d1) = ibu dari d1\\
		p  = relasi “anak dari” yaitu p(d1,d2) = d1 adalah anak dari d2
		
	\end{enumerate}
	Apakah arti ekspresi E berdasarkan interpretasi I dan interpretasi J ?
\end{frame}			

\section{Aturan Semantik}
\begin{frame}
	\frametitle{Aturan Semantik}
	Misal A adalah suatu ekspresi dan I adalah interpretasi untuk A yang meliputi domain tak kosong D. Maka nilai dibawah I ditentukan berdasarkan aturan semantik sebagai berikut :\\
	\begin{itemize}
		\item Nilai suatu konstanta a adalah elemen domain D.
		\item Nilai variabel x adalah elemen domain D.
		\item Nilai aplikasi $f_1$($t_1$, $t_2$, …, $t_n$) adalah elemen domain D dimana $f_1$($t_1$, $t_2$, …, $t_n$) f adalah fungsi yang diberikan kepada f dan $t_1$, $t_2$, …, $t_n$ adalah nilai term berdasarkan interpretasi I.
		\item Nilai Term kondisional if A then s else t adalah nilai term s jika A bernilai TRUE dan sama dengan nilai term t jika A bernilai FALSE.
		\item Nilai proposisi $p_1$($t_1$, $t_2$, …, $t_n$) adalah nilai kebenaran TRUE atau FALSE dimana p adalah relasi yang diberikan oleh interpretasi I dan nilai dari $t_1$, $t_2$, …, $t_n$ berdasarkan I.
		\item Aturan untuk penghubung logik (not, or, dsb) sama dengan aturan pada kalkulus proposisi.
	\end{itemize}	
\end{frame}	

\section{Interprestasi Yang Diperluas}
\begin{frame}
	\frametitle{Interprestasi Yang Diperluas}
	Misal I adalah suatu interpretasi yang mencakup domain D maka untuk sembarang variabel x dan elemen d pada domain D, interpretasi yang diperluas\\
	\begin{center}
		\textbf{$\langle x \leftarrow d \rangle \circ I $ }
	\end{center}
	adalah interpretasi yang mencakup domain D dimana :\\
	Variabel x diberikan nilai elemen pada domain D
	\begin{enumerate}
		\item Setiap variabel y (selain x) diberi nilai sama dengan elemen domain $y_1$ yaitu nilai berdasar interpretasi D. Jika y tidak memiliki nilai berdasar I maka y juga tidak memiliki nilai berdasar $\langle x \leftarrow d \rangle \circ I $.
		\item Setiap konstanta a, simbol fungsi f, dan simbol predikat p diberi nilai sesuai dengan nilai aslinya yaitu $a_I$, $f_I$, $p_I$.
	\end{enumerate}
	\textbf{Sifat interpretasi yang diperluas}\\	
	Jika I adalah interpretasi untuk kalimat berbentuk\\
	(FOR ALL x) A dan (FOR SOME x) A,\\
	maka $\langle x \leftarrow d \rangle \circ I $ adalah interpretasi yang berlaku untuk A juga
\end{frame}	

\begin{frame}
	\frametitle{Interprestasi Yang Diperluas}
	Contoh :\\
	\begin{enumerate}
		\item  I adalah interpretasi yang meliputi bilangan integer, dengan\\ 
		x = 1\\
		y = 2\\
		Maka perluasan interpretasi terhadap I :\\
		$\langle x \leftarrow 3 \rangle \circ I $\\
		akan memberikan nilai :\\
		x = 3\\
		y = 2
		\item I adalah interpretasi yang meliputi bilangan integer, dengan\\ 
		f adalah simbol fungsi biner\\
		+ adalah fungsi penambahan integer\\
		maka :\\
		$\langle f \leftarrow + \rangle \circ I $ adalah interpretasi yang meliputi domain bilangan integer dengan f fungsi penambahan +.
		
	\end{enumerate}
\end{frame}

\section{Aturan Semantik Untuk Kuantifier}
\begin{frame}
	\frametitle{Aturan Semantik Untuk Kuantifier}
	\textbf{Aturan FOR ALL}\\
	\begin{itemize}
		\item Kalimat (FOR ALL x) A bernilai TRUE berdasarkan interpretasi I jika:\\
		\underline{Untuk setiap elemen d} dari domain D menyebabkan subkalimat A bernilai TRUE berdasarkan interpretasi yang diperluas $\langle x \leftarrow d \rangle \circ I $
		\item Kalimat (FOR ALL x) A bernilai FALSE berdasarkan interpretasi I jika:\\
		\underline{Ada elemen d} dari domain D sedemikian sehingga subkalimat A bernilai FALSE berdasarkan interpretasi yang diperluas $\langle x \leftarrow d \rangle \circ I $
	\end{itemize}
	
\end{frame}

\begin{frame}
	\frametitle{Aturan Semantik Untuk Kuantifier}
	\textbf{Aturan FOR SOME}
	\begin{itemize}
		\item Kalimat (FOR SOME x) A bernilai FALSE berdasarkan interpretasi I jika:\\
		\underline{Untuk setiap elemen d} dari domain D menyebabkan subkalimat A bernilai FALSE berdasarkan interpretasi yang diperluas $\langle x \leftarrow d \rangle \circ I $
		\item Kalimat (FOR SOME x) A bernilai TRUE berdasarkan interpretasi I jika:\\
		\underline{Ada elemen d} dari domain D sedemikian sehingga subkalimat A bernilai TRUE berdasarkan interpretasi yang diperluas $\langle x \leftarrow d \rangle \circ I $
	\end{itemize}	
\end{frame}

\begin{frame}
	\frametitle{Aturan Semantik Untuk Kuantifier}
	Contoh :\\
	A : (FOR SOME x) p(x,y)
	Diberikan interpretasi I yang meliputi himpunan bilangan integer positif\\ 
	y = 2\\
	p : relasi “kurang dari”, yaitu $p_I$($d_1$,$d_2$) = $d_1<d_2$\\
	\vspace{5mm}	
	Berdasarkan aturan (FOR SOME x) maka\\
	(FOR SOME x) p(x, y) bernilai TRUE jika ada elemen dari D sehingga nilai p(x,y) bernilai TRUE berdasarkan interpretasi $\langle x \leftarrow d \rangle \circ I $\\
	\vspace{5mm}	
	Misal diambil d = 1 maka perluasan interpretasi menjadi sehingga berdasarkan aturan proposisi $\langle x \leftarrow 1 \rangle \circ I $ diperoleh bahwa\\
	p(1, 2) yaitu 1 $<$ 2 adalah TRUE		
\end{frame}

\begin{frame}
	\frametitle{Aturan Semantik Untuk Kuantifier}
	Contoh :\\
	B : IF (FOR ALL x) (FOR SOME y) p(x,y) THEN p(a,f(a))\\
	\vspace{5mm}
	 Misal I adalah interpretasi untuk B yang meliputi domain bilangan real positif dimana:\\
	 a = 1\\
	 f : fungsi “akar dari” yaitu f1(d) = $\sqrt{d}$\\
	 p : relasi “tidak sama dengan”, yaitu p1($d_1$,$d_2$) = $d_1 \neq d_2$\\
	 \vspace{5mm}
	 Misal diasumsikan bahwa B bernilai FALSE\\
	 Maka harus diperhatikan bahwa :\\ Antisenden : (FOR ALL x) (FOR SOME y) p(x,y) bernilai TRUE\\
	 Konsekuen : p(a,f(a)) bernilai FALSE
\end{frame}

\begin{frame}
	\frametitle{Aturan Semantik Untuk Kuantifier}
    Untuk lebih mudahnya, dimulai dari \underline{Konsekuen} karena bentuknya lebih sederhana. Berdasarkan aturan proposisi, maka nilai konsekuen p(a,f(a)) yaitu $1 \neq \sqrt{1}$ adalah 
    FALSE berdasarkan I\\
    \vspace{1mm}
    \underline{Antisenden} :\\
    berdasarkan Aturan (FOR ALL x)\\
    Untuk setiap elemen $d_1$ dari D, subkalimat (for some y) p(x,y) bernilai TRUE berdasarkan $\langle x \leftarrow d \rangle \circ I $\\
    \vspace{0.5mm}
    Berdasarkan Aturan (FOR SOME y)\\
    Untuk setiap elemen $d_1$ dari D, ada elemen $d_2$ sedemikian sehingga p(x,y) bernilai TRUE berdasarkan $\langle y \leftarrow d_2 \rangle \circ \langle x \leftarrow d_1 \rangle \circ I $\\
    \vspace{1mm}
    Misal ambil sembarang elemen domain dan $d_2 = d_1 + 1$\\
    Maka berdasarkan aturan proposisi, nilai p(x,y) yaitu p(d1, d2)\\
    Berarti p(d1, d1+1) menyatakan bahwa $d_1 \neq d_1 + 1$ adalah TRUE\\
    berdasarkan $\langle y \leftarrow d_2 \rangle \circ \langle x \leftarrow d_1 \rangle \circ I $\\
    \vspace{0.5mm}
    Jadi dapat disimpulkan bahwa kalimat B bernilai FALSE berdasarkan I.   
    
\end{frame}

\section{Kecocokan}
\begin{frame}
  \frametitle{Kecocokan}
  \begin{itemize}
  	\item Dua interpretasi dikatakan cocok jika keduanya memberi nilai yang sama untuk simbol-simbolnya \underline{atau} tidak memberi nilai untuk simbol-simbolnya keduanya. 
  	
  	\item Dua interpretasi I dan J cocok untuk ekspresi A jika nilai A berdasarkan I sama dengan nilai A berdasarkan J \underline{atau}
  	I dan J bukan interpretasi untuk A
  \end{itemize}
\end{frame}  

\begin{frame}
	\frametitle{Kecocokan}
	Contoh :\\
	Misalkan I adalah interpretasi yang meliputi bilangan integer dengan :\\
	a $\rightarrow$ 0\\ 
	b $\rightarrow$ 2\\
	x $\rightarrow$ -1\\ 
	f $\rightarrow$ fungsi suksessor $f_1$(d) = d + 1\\
	dan interpretasi J yang meliputi integer dengan :\\
	a $\rightarrow$ 0\\
	x $\rightarrow$ 1\\
	f $\rightarrow$	fungsi predesesor $f_1$(d) = d – 1   
	
\end{frame}

\begin{frame}
	\frametitle{Kecocokan}
	 \begin{itemize}
		\item I dan J cocok untuk konstanta a
		\item I dan J cocok untuk simbol predikat p
		\item I dan J tidak cocok untuk variabel x
		\item I dan J cocok untuk ekspresi f(x)
		\item I dan J cocok untuk ekspresi f(y)
		\item I dan J tidak cocok untuk ekspresi f(b), karena I adalah interpretasi untuk f(b) tetapi tidak untuk J 
	\end{itemize}
\end{frame}	

\section{Validitas}
\begin{frame}
	\frametitle{Validitas}
	Validitas di dalam kalkulus predikat didefinisikan hanya untuk kalimat tertutup, yaitu kalimat yang tidak memiliki variabel bebas.\\
	\vspace{2mm}
	\textbf{Definisi}\\
	Sebuah kalimat A dikatakan \underline{valid} jika kalimat tersebut bernilai TRUE berdasarkan setiap interprestasi untuk A.\\
	\vspace{2mm}
	Pembuktian validitas kalimat dapat menggunakan :\\
	\begin{itemize}
		\item Dengan membuktikan bahwa kalimat tertutup A adalah VALID	(biasanya lebih “enak” untuk kalimat-kalimat yang memiliki penghubung logik : 
		IFF, AND, NOT)
		\item Dengan membuktikan bahwa kalimat tertutup A adalah TIDAK VALID dengan	cara mencari satu interpretasi tertentu yang menyebabkan kalimat tersebut bernilai FALSE. 
		(biasanya untuk kalimat-kalimat yang memiliki penghubung logik : IF-THEN, OR)
	\end{itemize}
	
\end{frame}	

\begin{frame}
	\frametitle{Pembuktian Validitas}
	Misalkan ingin dibuktikan validitas kalimat A berikut :\\
	A : [ NOT (FOR ALL x) p(x) ] IFF [ (FOR SOME x) NOT p(x) ]\\
	\vspace{2mm}
	Berdasarkan aturan IFF, cukup diperlihatkan bahwa :\\
	NOT (FOR ALL x) p(x) ] dan [ (FOR SOME x) NOT p(x) ] memiliki nilai kebenaran yang sama berdasarkan setiap interpretasi, atau dengan kata lain subkalimat pertama bernilai TRUE tepat bila subkalimat kedua juga bernilai TRUE. 
\end{frame}	

\begin{frame}
	\frametitle{Pembuktian Validitas}
	A : [ NOT (FOR ALL x) p(x) ] IFF [ (FOR SOME x) NOT p(x) ]\\
	\vspace{2mm}
	Misalkan terdapat sebarang interpretasi I untuk A, maka\\
	NOT (FOR ALL x) p(x)  bernilai TRUE berdasarkan I\\
	Tepat bila (berdasarkan aturan NOT)\\
	\hspace{5mm} (FOR ALL x) p(x) bernilai FALSE berdasarkan I\\
	Tepat bila berdasarkan (FOR ALL x)\\
	\hspace{5mm} Ada elemen d di dalam domain D\\
	Sehingga p(x) bernilai FALSE berdasarkan $\langle x \leftarrow d \rangle \circ I $\\	
	Tepat bila berdasarkan aturan NOT\\
	\hspace{5mm} Ada elemen d di dalam domain D\\
	sehingga NOT p(x) bernilai TRUE berdasarkan $\langle x \leftarrow d \rangle \circ I $\\ 
	Tepat bila berdasarkan aturan (FOR SOME x)\\
	\hspace{5mm} (FOR SOME x) NOT p(x) bernilai TRUE berdasarkan Interpretasi I
\end{frame}

\begin{frame}
	\frametitle{Pembuktian Validitas}
	Misalkan ingin dibuktikan validitas kalimat B berikut :\\
	B :	IF (FOR SOME y) (FOR ALL x) q(x, y) THEN\\ 
	\hspace{5mm}(FOR ALL x) (FOR SOME y) q(x y)\\
	\vspace{2mm}
	Asumsikan bahwa B tidak valid, sehingga bahwa untuk suatu interpretasi I untuk B.\\
	\vspace{2mm}
	Jika Antisenden : (FOR SOME y) (FOR ALL x) q(x, y)\\
	\hspace{30mm}bernilai TRUE berdasarkan I\\
	maka konsekuen : (FOR ALL x) (FOR SOME y) q(x, y)\\
	\hspace{30mm}bernilai FALSE berdasarkan I
	
\end{frame}

\begin{frame}
	\frametitle{Pembuktian Validitas}
	Karena Antisenden bernilai TRUE berdasarkan  I,\\
	maka (berdasarkan aturan FOR SOME y)\\
	Ada elemen $d_1$ di dalam domain D sehingga (FOR ALL x) q(x, y) bernilai TRUE 
	berdasarkan $\langle x \leftarrow d_1 \rangle \circ I $\\
	\vspace{1mm}
	Tepat bila berdasarkan aturan FOR ALL x\\
	Ada elemen $d_1$ di dalam domain D sedemikian sehingga untuk setiap elemen $d_2$ di dalam domain D sedemikian sehingga q(x, y) bernilai TRUE berdasarkan $\langle x \leftarrow d_2 \rangle \circ \langle y \leftarrow d_1 \rangle \circ I $ ......... (1)\\
	\vspace{1mm}
	Karena konsekuen bernilai FALSE berdasarkan I,\\ 
	Maka (berdasarkan aturan FOR ALL x)Ada elemen $e_1$ di dalam domain D sehingga (FOR SOME y) q(x, y) bernilai FALSE 
	berdasarkan $\langle x \leftarrow e \rangle \circ I $\\
	\vspace{1mm}
	Tepat bila (berdasarkan aturan FOR SOME y)\\
	Ada elemen $e_1$ di dalam domain D sedemikian sehingga untuk semua elemen $e_2$ di dalam domain D sedemikian sehingga q(x, y) bernilai FALSE berdasarkan $\langle y \leftarrow e_2 \rangle \circ \langle x \leftarrow e_1 \rangle \circ I $ ......... (2)
	
\end{frame}

\begin{frame}
	\frametitle{Pembuktian Validitas}
	Berdasarkan (1) dan (2) kita dapat mengambil nilai elemen $d_1$ sama dengan dan $d_2$ sama dengan $e_1$, sehingga dari (1) diperoleh :\\
	\vspace{2mm}
	q(x y) bernilai TRUE berdasarkan $\langle x \leftarrow e_1 \rangle \circ \langle y \leftarrow d_1 \rangle \circ I $ ......... (3)\\
	dan dari (2) diperoleh\\
	 q(x, y) bernilai FALSE berdasarkan $\langle y \leftarrow d_1 \rangle \circ \langle x \leftarrow e_1 \rangle \circ I $ ......... (4)\\
	 Karena variabel x dan y berbeda, maka interpretasi\\
	 $\langle x \leftarrow e_1 \rangle \circ \langle y \leftarrow d_1 \rangle \circ I $ dan $\langle y \leftarrow d_1 \rangle \circ \langle x \leftarrow e_1 \rangle \circ I $\\
	 adalah identik, sehingga terlihat bahwa (3) dan (4) saling berkontradiksi.\\
	 \vspace{5mm}
	 \textbf{Berarti asumsi bahwa B tidak valid adalah tidak benar, sehingga B \underline{VALID}} 
\end{frame}									

\end{document}